%==============================================================================
\section{Affinity Propagation}
\label{sec:affinity-propagation}
%==============================================================================

\begin{dinhnghia}[title={Affinity Propagation}]
Thuật toán xác định các \textbf{exemplar} (điểm đại diện) và gán mỗi mẫu vào một exemplar.
\textbf{Không cần} chỉ định trước số cụm --- hữu ích cho phân tích khám phá (Frey \& Dueck, 2007).
\end{dinhnghia}

\subsection{Responsibility và Availability}

\begin{phuongphap}[title={Hai ma trận lặp}]
\begin{itemize}
  \item \textbf{Responsibility} $R(i,k)$: mức phù hợp của điểm $k$ làm exemplar cho điểm $i$.
  \item \textbf{Availability} $A(i,k)$: mức điểm $i$ ``muốn'' chọn $k$ làm exemplar.
\end{itemize}
Thuật toán \textbf{luân phiên} cập nhật hai ma trận đến hội tụ; exemplar cuối chọn theo giá trị lớn nhất trong ma trận responsibility.
\end{phuongphap}

\subsection{Triển khai Python (scikit-learn)}

\begin{vidu}[title={Ví dụ AffinityPropagation}]
\begin{lstlisting}
from sklearn.cluster import AffinityPropagation
from sklearn.datasets import make_blobs
import matplotlib.pyplot as plt

X, _ = make_blobs(n_samples=100, centers=4, random_state=0)

af = AffinityPropagation(preference=-50)
af.fit(X)

print("Cluster labels:", af.labels_)
print("Exemplars:", af.cluster_centers_indices_)

plt.figure(figsize=(7.5, 3.5))
plt.scatter(X[:, 0], X[:, 1], c=af.labels_, cmap='viridis')
plt.scatter(af.cluster_centers_[:, 0], af.cluster_centers_[:, 1],
            marker='x', color='red')
plt.show()
\end{lstlisting}
\end{vidu}

\begin{vidu}[title={Output mẫu trên terminal}]
\begin{verbatim}
Cluster labels: [3 0 3 3 3 3 1 0 0 0 0 0 0 0 0 2 3 3 1 2 2 0 1 2 3 1 3 3 2 2

 2 0 2 2 1 3 0 2 0 1 3 1 0 1 1 0 2 1 3 1 3 2 1 1 1 0 0 2 2 0 0 2 2 3 2 0 1 1 2

 3 0 2 3 0 3 3 3 1 2 2 2 0 1 1 2 1 2 2 3 3 3 1 1 1 1 0 0 1 0 1]

Exemplars: [9 41 51 74]
\end{verbatim}
\end{vidu}

\begin{ynghia}[title={Preference và Damping}]
  \begin{itemize}
    \item \textbf{Preference}: càng cao càng nhiều exemplar; càng thấp càng ít cụm.
    \item \textbf{Damping}: hệ số làm mượt hội tụ; damping lớn $\Rightarrow$ hội tụ chậm hơn.
  \end{itemize}
\end{ynghia}

\tsimg{12_affinity_propagation/img01.jpg}

\subsection{Ưu và nhược điểm}

\begin{tinhchat}[title={Ưu điểm}]
\begin{itemize}
  \item Tự xác định \textbf{số cụm}.
  \item Cụm \textbf{hình dạng/kích thước tùy ý}.
  \item Chịu được dữ liệu \textbf{nhiễu hoặc thiếu}.
  \item Tương đối \textbf{ổn định} với khởi tạo.
\end{itemize}
\end{tinhchat}

\begin{luuy}[title={Nhược điểm}]
\begin{itemize}
  \item \textbf{Tốn tài nguyên} trên dataset lớn hoặc nhiều feature.
  \item Có thể hội tụ \textbf{chưa tối ưu} khi dữ liệu biến thiên/nhiễu cao.
  \item \textbf{Nhạy} \texttt{damping}.
  \item Có thể sinh nhiều cụm \textbf{quá nhỏ} khó diễn giải.
\end{itemize}
\end{luuy}
